% !PRINTED 130
\weberpage{130}
Setzen wir
\weq{18}{
E\!\left(\begin{matrix}\alpha, & \beta \\ \gamma, & \delta\end{matrix}, \omega\right) = \varepsilon\sqrt{\alpha+\beta\omega},
}
so ist $\varepsilon$ eine 24ste Einheitswurzel, deren Produkt mit $\sqrt{\alpha+\beta\omega}$
durch (15) vollständig bestimmt ist, und es ergibt sich die Transformation
 der $\eta$-Funktion
\weq{19}{
\eta\!\left(\frac{\gamma+\delta\omega}{\alpha+\beta\omega}\right) = \varepsilon\sqrt{\alpha+\beta\omega}\ \eta(\omega).
}
Für $\varepsilon^{12}$ findet man
\weq{20}{
\varepsilon^{12} = (-1)^{\alpha\beta+\gamma\delta+\beta\gamma}.
}

\begin{center}
\textbf{§ 39. \gesperrt{Lineare Transformation der $\vartheta$-Funktionen.}}
\end{center}

Die Transformationsformeln der $\vartheta$-Funktionen sind Folgen
der Grundeigenschaften der $\sigma$-Funktion, durch die Substitution
\[
(\omega_1', \omega_2') = \begin{pmatrix}\alpha, & \beta \\ \gamma, & \delta\end{pmatrix}(\omega_1, \omega_2)
\]
ungeändert zu bleiben. Aus § 37, (4) ergibt sich hiernach, wenn
man $\omega_1, \omega_2, \eta_1$ durch $\omega_1', \omega_2', \eta_1'$ ersetzt:
\[
\frac{\omega_1\, e^{\frac{\eta_1 u^2}{\omega_1}}\, \vartheta_{11}\!\left(\dfrac{u}{\omega_1}, \dfrac{\omega_2}{\omega_1}\right)}{\vartheta'_{11}\!\left(0, \dfrac{\omega_2}{\omega_1}\right)}
=
\frac{\omega_1'\, e^{\frac{\eta_1' u^2}{\omega_1'}}\, \vartheta_{11}\!\left(\dfrac{u}{\omega_1'}, \dfrac{\omega_2'}{\omega_1'}\right)}{\vartheta'_{11}\!\left(0, \dfrac{\omega_2'}{\omega_1'}\right)},
\]
oder, weil nach § 35, (10), (13), (14)
\weq{1}{
\frac{\eta_1}{\omega_1} - \frac{\eta_1'}{\omega_1'}
= \frac{\omega_1'\eta_1 - \eta_1'\omega_1}{\omega_1\omega_1'}
= \beta\,\frac{\omega_2\eta_1 - \omega_1\eta_2}{\omega_1\omega_1'}
= \frac{\pi i\beta}{\omega_1\omega_1'}
}
ist, wenn wir:
\weq{2}{
\begin{aligned}
v &= \frac{u}{\omega_1}, & \omega &= \frac{\omega_2}{\omega_1},\\
v' &= \frac{u}{\omega_1'} = \frac{v}{\alpha+\beta\omega}, &
\omega' &= \frac{\omega_2'}{\omega_1'} = \frac{\gamma+\delta\omega}{\alpha+\beta\omega}
\end{aligned}
}
setzen:
\[
\frac{\vartheta_{11}(v', \omega')}{\vartheta'_{11}(0, \omega')}
= \frac{e^{\pi i\beta v v'}}{\alpha+\beta\omega}\ \frac{\vartheta_{11}(v, \omega)}{\vartheta'_{11}(0, \omega)}.
\]
Es ist ferner
\[
\vartheta'_{11}(0, \omega) = 2\pi\,\eta(\omega)^3, \qquad
\vartheta'_{11}(0, \omega') = 2\pi\,\eta(\omega')^3,
\]
woraus nach § 38, (19):
\[
\vartheta'_{11}(0, \omega') = \varepsilon^3\sqrt{\alpha+\beta\omega}^{\,3}\ \vartheta'_{11}(0, \omega),
\]
%CONT>>
